% =====================================================================
%  Currículo — versão ATS (Applicant Tracking System)
%  André Solano Ferreira Rodrigues Maiolini
%  Compilar com:  xelatex cv_maiolini_ats.tex
%
%  Diretrizes aplicadas:
%   - coluna única, fluxo de leitura linear (sem paracol, sem minipage)
%   - sem tikz, sem imagens, sem cor, sem caixas de texto
%   - títulos de seção com nomenclatura padrão de mercado
%   - contatos em linhas rotuladas, um por linha
%   - sem modo matemático, sem entreletras artificial, sem ligaduras exóticas
%   - fonte Carlito (métrica de Calibri), com mapeamento Unicode correto
% =====================================================================
\documentclass[10pt,a4paper]{article}

\usepackage{fontspec}
\usepackage[a4paper,top=1.3cm,bottom=0.9cm,left=1.5cm,right=1.5cm]{geometry}
\usepackage{titlesec}
\usepackage{enumitem}
\usepackage[hidelinks]{hyperref}
\usepackage{microtype}

\setmainfont{Carlito}[
  BoldFont={Carlito Bold}, ItalicFont={Carlito Italic}]

\setlength{\parindent}{0pt}
\raggedright
\linespread{1.0}
\pagestyle{empty}

\hypersetup{
  pdftitle={Currículo - André Solano Ferreira Rodrigues Maiolini},
  pdfauthor={André Solano Ferreira Rodrigues Maiolini},
  pdfsubject={Controle e Sistemas Dinâmicos},
  pdfkeywords={sistemas embarcados, controle, FPGA, VHDL, SystemVerilog,
               elementos finitos, MATLAB, Simulink, C, C++, Python, RISC-V, FreeRTOS,
               arranjo sistólico, quantização int8}
}

% ---------- seções -------------------------------------------------
\titleformat{\section}
  {\normalsize\bfseries\MakeUppercase}{}{0pt}{}[\vspace{2pt}\hrule height 0.6pt]
\titlespacing*{\section}{0pt}{8pt}{5pt}

% ---------- listas --------------------------------------------------
\setlist[itemize]{leftmargin=13pt, labelsep=5pt, label=\textbullet,
  topsep=2pt, itemsep=2pt, parsep=0pt}

% ---------- cabeçalho de entrada -----------------------------------
% #1 cargo/título  #2 organização  #3 período e local
\newcommand{\entrada}[3]{%
  \textbf{#1}\par
  #2\par
  #3\par
  \vspace{2pt}}

\newcommand{\espaco}{\vspace{7pt}}

% =====================================================================
\begin{document}

% ---------- IDENTIFICAÇÃO E CONTATO ---------------------------------
{\LARGE\bfseries André Solano Ferreira Rodrigues Maiolini}\par
\vspace{3pt}
{\large Engenheiro | Controle e Sistemas Dinâmicos}\par
\vspace{7pt}

Telefone: (11) 97120-8137\par
E-mail: \href{mailto:andre.maiolini@hotmail.com}{andre.maiolini@hotmail.com}\par
LinkedIn: \href{https://linkedin.com/in/andremaiolini}{linkedin.com/in/andremaiolini}\par
GitHub: \href{https://github.com/Nyfeu}{github.com/Nyfeu}\par
Localização: Mooca, São Paulo - SP, Brasil\par

% ---------- RESUMO --------------------------------------------------
\section{Resumo Profissional}

Engenheiro Mecânico pelo Instituto Mauá de Tecnologia e formando em Engenharia
de Computação pela mesma instituição. Atuação entre a modelagem de sistemas
físicos - com ênfase em elementos finitos, dinâmica e vibrações - e o
desenvolvimento de hardware digital e sistemas embarcados. Iniciação
científica, estágio e trabalho de conclusão seguem a mesma abordagem:
confrontar o modelo com resultados experimentais. Busco atuar no
desenvolvimento de sistemas de controle e sistemas dinâmicos, aproximando
modelagem, computação e implementação embarcada em aplicações de tempo real.

% ---------- COMPETÊNCIAS --------------------------------------------
\section{Competências Técnicas}

\textbf{Controle e sistemas dinâmicos:} malhas de controle embarcadas, filtros
digitais, sistemas de tempo real, dinâmica e vibrações.\par
\vspace{3pt}
\textbf{Modelagem numérica e experimental:} Método dos Elementos Finitos
(FEM/FEA), Ansys Mechanical, Ansys Fluent (CFD), MSC Simufact, planejamento de
experimentos (DoE), ANOVA, análise modal.\par
\vspace{3pt}
\textbf{Hardware digital e sistemas embarcados:} VHDL, SystemVerilog, FPGA
(Xilinx Vivado, Nexys A7), ESP32, MSP430, FreeRTOS, RISC-V, UART, DMA,
arranjo sistólico, quantização int8, Hardware-in-the-Loop (HIL).\par
\vspace{3pt}
\textbf{Programação e ferramentas:} MATLAB, Simulink, C, C++, Python, Rust, Flutter, Git,
Linux, integração contínua (CI).\par

% ---------- IDIOMAS -------------------------------------------------
\section{Idiomas}

Português: nativo.\par
Inglês: intermediário.\par

% ---------- EXPERIÊNCIA ---------------------------------------------
\section{Experiência Profissional}

\entrada{Assistente de Coordenação}
        {Instituto Mauá de Tecnologia}
        {02/2026 - atual | São Caetano do Sul, SP}
\begin{itemize}
  \item Apoio técnico e didático em Arquitetura e Organização de Computadores e
        orientação de projetos de graduação, incluindo desenvolvimento de
        firmware e hardware para dispositivos IoT.
\end{itemize}

\espaco
\entrada{Estagiário de Engenharia de Sistemas Embarcados}
        {Intelicity}
        {02/2025 - 01/2026 | São Paulo, SP}
\begin{itemize}
  \item Desenvolvimento integral de plataforma automatizada de posicionamento
        angular em dois eixos (CAD, fabricação aditiva, eletrônica, firmware e
        aplicativo de operação), com posicionamento em malha fechada por
        encoders magnéticos absolutos. Empregada como referência angular na
        calibração de acelerômetros MEMS, com resultados publicados no
        COBRAE/GEOSUL 2025.
  \item Firmware em C para ESP32 e MSP430 sobre FreeRTOS, com filtros digitais
        aplicados à cadeia de aquisição de sinais dos sensores.
  \item Dashboards em Python para bancada de testes e aplicativo em Flutter
        para operação.
\end{itemize}

\espaco
\entrada{Aluno Pesquisador - Iniciação Científica}
        {Instituto Mauá de Tecnologia}
        {02/2022 - 01/2024 | São Caetano do Sul, SP}
\begin{itemize}
  \item Modelagem termomecânica da fusão seletiva a laser (SLM/PBF-L) de
        componentes em Ti6Al4V por Elementos Finitos, utilizando o MSC
        Simufact, sob planejamento fatorial e ANOVA, com validação experimental
        inferior a 4,5\% para distorções térmicas e a 1\% para tensões
        residuais medidas por difratometria de raios-X.
  \item Caracterização experimental de ligas refratárias (WC e NbC com Co e Ni)
        para manufatura aditiva por PBF-L, incluindo análises por Microscopia
        Eletrônica de Varredura (MEV) e levantamento experimental de curvas de
        escoabilidade a seco.
\end{itemize}

\espaco
\entrada{Monitor e Professor Particular}
        {Instituto Mauá de Tecnologia e autônomo}
        {02/2020 - 12/2023 | São Paulo, SP}
\begin{itemize}
  \item Aulas particulares de Mecânica Vibratória e Controle, Resistência dos
        Materiais e Cálculo Diferencial e Integral I e II.
  \item Monitorias de Resistência dos Materiais e de Cálculo I, com elaboração
        de materiais didáticos complementares.
\end{itemize}

% ---------- FORMAÇÃO ------------------------------------------------
\section{Formação Acadêmica}

\entrada{Bacharelado em Engenharia de Computação}
        {Instituto Mauá de Tecnologia}
        {02/2024 - 12/2026 (em curso)}
\begin{itemize}
  \item Bolsa Melhor Aluno 2025 (IMT, 2026): maior Coeficiente de Rendimento
        entre os alunos da quarta série de Engenharia, com bolsa integral de
        100\%.
\end{itemize}

\espaco
\entrada{Bacharelado em Engenharia Mecânica}
        {Instituto Mauá de Tecnologia}
        {02/2019 - 12/2023 (concluído)}
\begin{itemize}
  \item Melhor Aluno da Turma de 2023 (IMT, 2024) e Prêmio Instituto de
        Engenharia de Melhor Aluno da Turma de Formandos de 2023.
  \item 27º Prêmio CREA-SP de Formação Profissional (CREA-SP, 2024).
\end{itemize}

% ---------- PROJETOS ------------------------------------------------
\section{Projetos Técnicos}

\entrada{SuperNova-RV: exploração microarquitetural de núcleos RISC-V}
        {Projeto pessoal}
        {2026 - em andamento}
\begin{itemize}
  \item Núcleo RV32I em SystemVerilog, de single-cycle a pipeline de cinco
        estágios com tratamento de hazards, com métricas de CPI/IPC e
        verificação automatizada em CI.
\end{itemize}

\espaco
\entrada{Plataforma didática: SoC RISC-V com aceleração neural em FPGA}
        {Trabalho de Conclusão de Curso - Instituto Mauá de Tecnologia}
        {2025 - 2026}
\begin{itemize}
  \item Coprojeto de hardware/software de um sistema-em-chip (SoC) em
        VHDL-2008, validado em FPGA Nexys A7: núcleo RV32I multiciclo com
        Zicsr, NPU em arranjo sistólico 4x4, DMA e interconexão crossbar.
  \item Mitigação do gargalo de memória na NPU (pesos residentes, im2col em
        hardware, streaming, sequenciamento de tiles, modo GEMV, fusão de
        camadas e descritores): 12,2 MACs úteis por ciclo, com o arranjo em
        computação 76\% do tempo.
  \item Inferência de CNN int8 (MNIST) em 1.506 ciclos, com perda de 0,13 ponto
        percentual de acurácia na quantização.
\end{itemize}

\espaco
\entrada{Otimização de rotor de misturador de alto cisalhamento}
        {Parceria SEMCO e Instituto Mauá de Tecnologia}
        {02/2023 - 06/2023}
\begin{itemize}
  \item Primeira colocação no desafio de otimização da geometria do rotor.
        Estudo paramétrico com cerca de 20 simulações em Ansys Fluent, com
        múltiplos referenciais móveis (MRF) e turbulência k-epsilon realizable,
        variando número, inclinação e altura das pás.
  \item Verificação da geometria selecionada por interação fluido-estrutura e
        análise modal: excitação de 56,7 Hz do rotor confrontada com os modos
        naturais da estrutura.
\end{itemize}

% ---------- PUBLICAÇÕES ---------------------------------------------
\section{Publicações}

\begin{itemize}
  \item \textbf{Influência dos nanofluidos na tensão residual do aço AISI 316L
        na manufatura híbrida.} CONEM 2024, ABCM. DOI:
        10.26678/abcm.conem2024.con24-0187.
  \item \textbf{Numerical and Experimental Analysis of the Influence of
        Manufacturing Parameters in Additive Manufacturing SLM-PBF on Residual
        Stress and Thermal Distortion in Parts of Titanium Alloy Ti6Al4V.}
        Springer, Lecture Notes in Mechanical Engineering, 2024. DOI:
        10.1007/978-3-031-58006-2\_15.
  \item \textbf{Towards the Flowability and Spreadability of Cemented
        Carbides and Cermet Powders for Additive Manufacturing: Experimental
        and Numerical Approach - Part 1.} COBEM 2023, ABCM. DOI:
        10.26678/ABCM.COBEM2023.COB2023-0100.
  \item \textbf{Análise numérica e experimental da influência dos parâmetros de
        fabricação em manufatura aditiva SLM-PBF nas tensões residuais e
        distorções térmicas em peças de Ti6Al4V.} COBEF 2023, ABCM. DOI:
        10.26678/ABCM.COBEF2023.COF23-0113.
  \item \textbf{Uso de estratégias ativas no ensino de Método dos Elementos
        Finitos no curso de Engenharia Mecânica.} COBENGE 2022, ABENGE. DOI:
        10.37702/2175-957X.COBENGE.2022.3795.
\end{itemize}

% ---------- CERTIFICAÇÕES -------------------------------------------
\section{Certificações}

\begin{itemize}
  \item Control System Design with MATLAB and Simulink - MathWorks, 2026.
  \item Academia ESP32 Profissional - Embarcados, 2026.
  \item Análise Dinâmica com Ansys - Instituto ESSS, 2024.
  \item Análise Não-Linear com Ansys - Instituto ESSS, 2023.
  \item Formação Linguagem C e C++ - Alura, 2023.
  \item Formação Linguagem Rust - Alura, 2023.
  \item MATLAB: do básico ao avançado - Udemy, 2022.
\end{itemize}

\end{document}